\documentclass[10pt,a4paper]{amsart}
\usepackage[margin=19mm]{geometry}
\usepackage{amsmath,amssymb,mathtools}
\usepackage[scr=rsfs,cal=euler]{mathalfa}
\usepackage{xfrac}
\usepackage[unicode,hidelinks,bookmarks=false]{hyperref}
\newcommand{\ie}{i.e.\ }
\newcommand{\cf}{cf.\ }
\newcommand{\resp}{respectively\ }
\newcommand{\Subsection}{\subsection}
\providecommand{\nobreakdash}{\nobreak\hbox{-}}
\setlength{\emergencystretch}{2em}
\multlinegap0pt
\usepackage{bm}
\usepackage{bbm}
\usepackage{dsfont}
\usepackage{xspace}
\usepackage{cancel}
\usepackage{mathtools}
\usepackage{amsthm}
\makeatletter
\@ifundefined{thm}{\newtheorem{thm}{Theorem}}{}
\makeatother
\newcommand{\G}{\mathcal{G}}
\newcommand{\trans}{^{\top}}
\title[The Inverse Symplectic Eigenvalue Problem of a Graph]{The Inverse Symplectic Eigenvalue Problem of a Graph}
\author{Himanshu Gupta, Leslie Hogben, Bryan Shader,, Tony Wong}
\date{Source: arXiv:2601.11912v1 (CC BY 4.0)}
\begin{document}
\maketitle

\noindent\emph{Source and licence.} The Inverse Symplectic Eigenvalue Problem of a Graph. Version arXiv:2601.11912v1 (CC BY 4.0), distributed under the Creative Commons Attribution 4.0 International licence, as recorded in the arXiv licence metadata for this exact version and shown on the abstract page https://arxiv.org/abs/2601.11912v1. Full rights notice: licenses/arxiv-2601.11912-CC-BY-4.0.txt.\par\smallskip

\noindent\textit{Editorial scope.} Counted unit: the Thm whose source label is \texttt{t:PDirred-sparseLB} in the article (source0.tex lines 778--783, source label \texttt{t:PDirred-sparseLB}), delivered together with the article's own complete original proof of it (source0.tex lines 785--988, one \texttt{proof} environment of 204 lines). No mathematics is written, repaired or re-derived: statement and proof are the article's own text line for line. Required context retained uncounted: none. Every definition, display and label of those lines is retained unchanged. Omitted from this package and not redistributed: 1--777, 784--784, 989--2799. Nothing inside a retained excerpt is omitted or altered: every retained body is delivered exactly as the article's lines are written, so no omission receipt of any kind accompanies this package. Citations: the retained excerpts carry no citation command at all, so no bibliography entry of the article is transcribed and no reference is redistributed. Reference adaptation: none was needed and none was made; every reference inside the delivered unit resolves inside it, so no unresolved reference or citation is produced. Printed slips kept literal: no printed word, symbol, hypothesis or number of the counted statement or of its proof was altered. Layout and class: the article's own class is \texttt{article} with packages including amsmath, amssymb, amsthm, amsfonts, xcolor, enumerate, hyperref, cleveref, graphicx, soul, url, dsfont, color, tikz; the wrapper is the lane's pinned \texttt{amsart} editorial wrapper at 10pt on A4, which restates the theorem-like environments the retained text needs and adds the article's own macro definitions that the retained text needs (\texttt{\textbackslash G}, \texttt{\textbackslash trans}), with extra wrapper packages (bm, bbm, dsfont, xspace, cancel, mathtools). The delivered numbering is the unit's own. Assets: no retained span contains a figure environment, a graphics-inclusion command, a PDF-page inclusion, a table or a TikZ picture; no image file is copied into this package, the package declares no asset dependency and no image file is redistributed with it. Licence: version arXiv:2601.11912v1 of this article is distributed under the Creative Commons Attribution 4.0 International licence, as recorded in the arXiv licence metadata for this exact version and shown on the abstract page https://arxiv.org/abs/2601.11912v1. A full rights notice is delivered at licenses/arxiv-2601.11912-CC-BY-4.0.txt. Characters: every retained excerpt is pure ASCII, so the delivered engine, XeLaTeX with its default Latin Modern text font, reports no missing character.
\begin{thm}\label{t:PDirred-sparseLB}
Let $M$ be an $n\times n$ irreducible, 
positive definite matrix. 
Then 
\[\#(M)+ \#(M^{-1}) \geq 8n-8.\]
\end{thm}

\begin{proof}
The proof is by induction on $n$.
{For $n=1$, $\#(M)+ \#(M^{-1})= 1+1 > 0=8\cdot 1 -8$.} 
For $n=2$, the irreducibility of $M$ implies 
that $\#(M) + \#(M^{-1})=4 +4 =8\cdot 2 -8$.

Now assume that $n>2$ and proceed by induction.
Throughout {this proof} we use $N$ to denote $M^{-1}$.  Since $M$ is positive definite, every diagonal entry is nonzero, so the number of nonzero vertices in row $i$ of $M$ is $1+\deg_{\G(M)}(i)$, and similarly for $N$.
If each vertex of $\G(M)$ and $\G(N)$ has degree at least 3, then 
\[ 
\#(M)+ \#(N) \geq 4n+4n=8n> 8n-8,
\]
as desired. Thus, we may assume that $\G(M)$ or $\G(N)$ 
has a vertex of degree at most $2$. We consider several cases. 

\bigskip\noindent
{\bf Case 1.} $\G(M)$ or $\G(N)$
has a vertex of degree $1$.

Without loss of generality, we can assume that  vertex
$1$ is a pendant vertex adjacent to  vertex $2$ in $\G(M)$.
Thus $M$ has the form 
\[
\begin{bmatrix}
m_{11} & m_{12} & 0 &\cdots & 0\\
m_{12} & m_{22} \\
0\\
\vdots\\
0 
\end{bmatrix}.
\]
 By considering the first row of the product  $MN=I$, we see that {$N_{\G(N)}[1]=N_{\G(N)}[2]$}. Since $\G(M)$, and hence $\G(N)$,
is connected and $n\geq 3$, vertex $1$ of $\G(N)$ has degree at least $2$. 

Let 
\begin{align*}
P&= \begin{bmatrix}
    1 & -m_{12}/m_{11}\\
    0 & 1
\end{bmatrix} \oplus I_{n-2},\\
\widehat{M}&= P\trans MP, \mbox{ and }\\
\widehat{N}&= P^{-1} N (P^{-1})\trans .
\end{align*}
Then 
$\widehat{M}$ and $\widehat{N}$ are
 positive definite matrices that are inverses of each other.  Examination of $\widehat{M}=P\trans MP$ shows that the  matrix $\widehat{M}$ has the form 
$[m_{11}] \oplus B$  where {(with the rows and columns of $B$ indexed by $2,\dots,n$)  $b_{22}=m_{22}-\frac{m_{12}^2}{m_{11}}$} and $b_{ij}=(M(1))_{ij}$ for $i,j\ge 2$ and $(i,j)\ne (2,2)$. Hence $\#(B)= \#(M(1))$.  Examination of $\widehat{N}=P^{-1} N (P^{-1})\trans$ shows that $\widehat{N}=
[1/m_{11}] \oplus N(1)$. 

Since $\G(M)$ is connected and vertex $1$ is a pendant vertex, $\G(M(1))$ is connected.
Hence $B$ is irreducible symmetric positive definite of order $(n-1)$ with inverse $N(1)$.
By induction, $\#(B) + \#(N(1))\geq 8(n-1)-8=8n-16$, implying that $\#(M(1))+ \#(N(1)) \geq 8n-16$.
The number of nonzeros in the first row and column of $M$ is 3, the number of nonzeros in the first row and column of $N$ is at least 5.
Therefore $\#(M)+ \#(N)\geq 8n-16 + 3 +5= 8n-8$,
as desired.


\bigskip\noindent 
{\bf Case 2.} All vertices of $\G(M)$ and $\G(N)$ have degree at least 2 and at least one has degree 2. 

Without loss of generality, we may assume that the neighbors of vertex $1$
in $\G(M)$ are vertices $2$ and $3$.  Thus $M$ has the form 
\[
M= \begin{bmatrix}
    m_{11} & m_{12} & m_{13} & 0 & \cdots & 0 \\
    m_{12}\\
    m_{13}
    \\
    0 \\
    \vdots\\
    0
\end{bmatrix}.
\]
Let 
\begin{align*}
P&= \begin{bmatrix}
    1 & -m_{12}/m_{11} & -m_{13}/m_{11}\\
    0 & 1 & 0 \\
    0 & 0 & 1
\end{bmatrix} \oplus I_{n-3},\\
\widehat{M}&= P\trans M P, \mbox{ and }\\
\widehat{N}&= P^{-1} N(P^{-1})\trans.
\end{align*}
Then $\widehat{M}$ and $\widehat{N}$ are symmetric positive definite matrices that are inverses of each other. Examination of $\widehat{M}=P\trans MP$ shows that  $\widehat{M}= [m_{11}] \oplus B$  where $b_{22}=-\frac{m_{12}^2}{m_{11}} + m_{22}$, $b_{23}=-\frac{m_{12}m_{13}}{m_{11}} + m_{23}$, $b_{33}=-\frac{m_{13}^2}{m_{11}} + m_{33}$, and $b_{ij}=m_{ij}$ for all other entries of $B$.
 Thus $\G(B)$ is  $\G(M(1))$ or is obtained from $\G(M(1))$ 
by deleting  or inserting  the edge joining vertices $2$ and $3$. Examination of $\widehat{N}=P^{-1} N (P^{-1})\trans$ shows that $\widehat{N}=
[1/m_{11}] \oplus N(1)$.  Since $B$ and $N(1)$ are inverses of each other, either   $\G(B)$ and $\G(N(1))$ are both connected or $\G(B)$ and $\G(N(1))$ are both disconnected.   

\bigskip\noindent
{\bf Subcase:} $\G(N(1))$ is connected.
\\
In this case, both  $\G(B)$  and $\G(N(1))$ are connected. So $B$ and $N(1)$ 
are irreducible, positive definite symmetric matrices of order $n-1$ 
and are inverses of each other. By induction, 
\[ \#(B) + \#(N(1))\geq 8(n-1)-8=8n-16.\]
The first row and column of $N$ contain at least $5$ nonzeros since all vertices of $\G(N)$ have degree at least $2$.
So $\#(N)\geq \#(N(1))+5$.  The first row and column of $M$ contain 
at least $5$ nonzeros, and $\#(M(1))\geq \#(B)-2$.  So, 
$\#(M) \geq \#(B)+3$.  It now follows that 
\[ \#(M) +\#(N) \geq 8n-16 + 5 + 3=8n-8,\] as desired. 

\bigskip\noindent
{\bf Subcase:} $\G(N(1))$ is disconnected.
\\
Since $B$ and $N(1)$ are inverses of each other, $B$ is also disconnected.  
Since $\G(B)$ and $\G(M(1))$ have the same edges except for possibly the edge joining $2$ and $3$, $\mathcal{G}(N(1))$ has exactly two 
connected components, one containing $2$ and 
the other containing $3$. Denote the vertex sets of the two connected components of $N(1)$ by $ \alpha$ and $ \beta$ where $2 \in \alpha$
and $3 \in \beta$.
Then  $M$ has the form 
\[ 
M= 
\left[ \begin{array}{ccc|c|c}
m_{11} & m_{12} & m_{13} & \mathbf 0^{\top} & \mathbf{0}^{\top}\\
m_{12} & m_{22} & m_{23} & \mathbf g^{\top} & \mathbf h^{\top} \\
m_{13} & m_{23} & m_{33} & \mathbf i^{\top} & \mathbf j^{\top}  \\ \hline
\mathbf 0 & \mathbf g &\mathbf i &  M[\alpha\setminus\{2\} ] & O 
\\ \hline 
\mathbf 0 & \mathbf h & \mathbf j &  O & M[\beta\setminus \{3\}] 
\end{array} \right] .  
\] 
and $N$ has the form 
\[
N= 
\left[ 
\begin{array}{ccc|c|c} 
n_{11} & n_{12} & n_{13} & \mathbf u^{\top} & \mathbf v^{\top} \\
n_{12} & n_{22}  & 0 & \mathbf x^{\top} & \mathbf 0^{\top}\\
n_{13} & 0 & n_{33} &  \mathbf 0^{\top} & \mathbf y^{\top} \\ \hline 
\mathbf u & \mathbf x &  \mathbf 0 &  N[\alpha\setminus\{2\}] & O \\ \hline 
\mathbf v & \mathbf {0} & {\mathbf y} & O & N[\beta
\setminus\{3\} ]
\end{array} 
\right] .
\] 
Note that  
\[ 
P^{\top}MP =
\left[ \begin{array}{ccc|c|c}
m_{11} & 0 & 0 & \mathbf 0^{\top} & \mathbf{0}^{\top}\\
0 & m_{22}-m_{12}^2/m_{11} & m_{23}-m_{12}m_{13}/m_{11} & \mathbf g^{\top} & \mathbf h^{\top} \\
0 & m_{23}-m_{12}m_{13}/m_{11} & m_{33}-m_{13}^2/m_{11}  & \mathbf i^{\top} & \mathbf j^{\top}  \\ \hline
\mathbf 0 & \mathbf g &\mathbf i &  M[\alpha\setminus \{2\} ] & O 
\\ \hline 
\mathbf 0 & \mathbf h & \mathbf j &  O & M[\beta\setminus \{ 3 \}] 
\end{array} \right] ,  
\] 
and 
$ P^{-1}N (P^{\top})^{-1} $
equals 
\begin{center}
\resizebox{.95\linewidth}{!}{$\left[ 
\begin{array}{ccc|c|c} 
 1/m_{11} & n_{12} + n_{22}m_{12}/m_{11} & n_{13} + n_{33}m_{13}/m_{11} & \mathbf u^{\top}+ (m_{12}/m_{11})\mathbf x^{\top}  & \mathbf v^{\top} + (m_{13}/m_{11}) \mathbf y^{\top} \\
n_{12}+n_{22}m_{12}/m_{11}  & n_{22} & 0 & \mathbf x^{\top} & \mathbf 0^{\top}\\
n_{13} +n_{33}m_{13}/m_{11} & 0 & n_{33} &  \mathbf 0^{\top} & \mathbf y^{\top} \\ \hline 
\mathbf u+ (m_{12}/m_{11})\mathbf x & \mathbf  x &  \mathbf 0 &  N[\alpha \setminus \{ 2\} ] & O \\ \hline 
\mathbf v + (m_{13}/m_{11})
\mathbf y & \mathbf {0} & \mathbf {y} & O & N[\beta\setminus \{3\} ]
\end{array} 
\right]$}.
\end{center}
Since $P^{\top} MP[\alpha , \beta]=O$, $m_{23}=m_{12}m_{13}/m_{11}$.  
Thus $m_{23} \neq 0$, that is,  there is an edge joining vertex $2$ and vertex $3$ in $\G(M)$. 
Since $P^{\top}MP$ and $P^{-1}N(P^{\top})^{-1}$ are inverses of each other, each off-diagonal entry in the first row of 
$P^{-1} N(P^{\top})^{-1}$\emph  is zero. 
Thus, 
\begin{align*}
n_{12}& =-n_{22}m_{12}/m_{11}\\
n_{13}&=-n_{33}m_{13}/m_{11}\\
\mathbf u& = -(m_{12}/m_{11}) \mathbf x \\ 
\mathbf v&= -(m_{13}/m_{11}) \mathbf y. 
\end{align*}
In particular, this implies for $j \in \alpha \setminus \{2\} $,  vertex  $j$ is a neighbor of vertex $2$ in $\G(N)$ if 
and only if $j$ is a neighbor of $1$ in $\G(N)$.
Since each vertex of $\G(N)$ has degree at least $2$, there is some $j \in \alpha \setminus \{2\} $ that  is a neighbor of vertex $2$. 
Hence vertex $1$ has at least two neighbors in $\alpha$. Similarly, vertex $1$ has at least two neighbors in $\beta$, and the degree $d$ 
of vertex $1$ in 
$\G(N)$ is at least $4$.


Up to permutation similarity, 
$B= B[\alpha]  \oplus B[\beta ]$
and $N(1)= N[\alpha] \oplus N[\beta]$. So $B[\alpha]$ and $B[\beta]$ are irreducible positive definite matrices of order $|\alpha|$ and $|\beta|$ respectively, and their inverses are $N[\alpha]$ and $N[\beta]$ respectively. 
Hence, by induction $\#(B[\alpha]) + \#(N[\alpha])\geq 8|\alpha| - 8$, and 
$\#(B[\beta]) + \#(N[\beta])\geq 8|\beta|-8$.


Note that $N$ has $1+ 2 {d}$ nonzero entries in its first row and column,  $M$ has $5$ nonzero entries in its first row and column, and $\G(B)$ is $\G(M(1))$ with exactly one edge removed.  
So $\#(M(1))=
\#(B)+2$. Hence, $\#(M) =\#(B[\alpha]) + \#(B[\beta])+7$.
It follows that 
\begin{align*} 
\#(M) +\#(N)  
&\geq \#(B[\alpha]) + \#(B[\beta]) +7 + \#(N[\alpha]) + \#(N[\beta]) + (1+ 2d)
\\
&= \#(B[\alpha]) + \#(N[\alpha])+ \#(B[\beta]) + \#(N[\beta]) +8 + 2d \\
&\geq 8|\alpha|-8+ 8|\beta|-8 + 8+ 2d \\
&=8(n-1) + 2d -8 
\\
&\geq 8n-8.
\end{align*}
\qedhere 
\end{proof}

\end{document}
